Detectors of AI-generated text routinely report accuracy at or above 99\% on benchmarks such as DAIGT V2 and HC3, but this in-distribution number says little about real-world reliability. We decompose reported accuracy into three components, cross-dataset generalization, dependence on surface artifacts, and robustness to cheap adversarial attack, using a five-model pipeline (three classical baselines and fine-tuned BERT and DeBERTa) at full corpus scale. A contamination audit finds HC3 carries substantial internal duplication (7.16\%) essentially absent from DAIGT V2 (0.01\%). Strict one-way cross-dataset transfer, evaluated across three seeds, shows substantial degradation from both datasets in both directions. An artifact-cleaning ablation, run both zero-shot and via full retraining, finds a real negative accuracy delta for HC3 but not DAIGT V2, consistent with HC3's stronger surface artifact, a whitespace-before-punctuation cue present in 88.7\% of human answers versus 0.28\% of ChatGPT answers. An adversarial evaluation across typo injection, homoglyph substitution, and back-translation then finds HC3 accuracy collapses to near-chance under a typo attack, as low as 0.36 F1 from an in-domain 0.998, while DAIGT V2 stays comparatively robust, three independent experiments converging on the same conclusion for HC3. These results support a dataset-dependent verdict, HC3's headline accuracy substantially overstates real-world reliability while DAIGT V2's does so only modestly, a distinction invisible from in-distribution accuracy alone.
